\chapter{The MQTT / Sparkplug~B contract}
\label{ch:contract}

This chapter is the versioned integration contract of record between \prog and Ignition:
the topic grammar, the metric definitions, quality/freshness semantics, and the
lifecycle (NBIRTH/NDATA/NDEATH).

\begin{note}
Everything in this chapter describes behaviour that is \emph{implemented and tested} as of
Epic~1. Where the bridge does not yet do something the Sparkplug specification allows, it
is listed in section~\ref{sec:contract-limits} rather than described as though it worked.
\end{note}

\section{Contract version}

The node BIRTH carries a metric \code{Contract/Version} (\code{Int64}), currently
\textbf{1}. It is incremented on any breaking change to the topic grammar, to a metric
name, or to a metric unit. Publishing it means a consumer can \emph{see} the contract
change instead of having to notice it from broken tags.

\section{Topic grammar}

The namespace is \code{spBv1.0}. One EON node represents the bridge; each meter is a
\emph{device} beneath it, identified by its immutable serial number.

\begin{center}
\begin{tabular}{@{}l>{\ttfamily}l@{}}
\toprule
\normalfont\textbf{Level} & \normalfont\textbf{Topic} \\
\midrule
Node   & spBv1.0/\{group\}/\{TYPE\}/\{node\} \\
Device & spBv1.0/\{group\}/\{TYPE\}/\{node\}/\{serial\} \\
\bottomrule
\end{tabular}
\end{center}

\code{group} and \code{node} are configured; \code{serial} is the meter's serial number,
never a friendly name. Binding the topic to the serial is what makes it impossible to
attribute a value to the wrong meter after a device is renamed or replaced.

Identifiers are validated \emph{before} the bridge connects. A serial that cannot be a
topic level makes the bridge refuse to start, because the alternative — connecting,
never birthing, and publishing nothing — looks identical to a healthy but idle bridge.

\subsection{Message types emitted}

\begin{center}
\begin{tabular}{@{}>{\ttfamily}l l p{7.4cm}@{}}
\toprule
\normalfont\textbf{Type} & \textbf{Level} & \textbf{When} \\
\midrule
NBIRTH & node   & On every successful connect, including reconnects. \\
DBIRTH & device & Immediately after NBIRTH, one per meter. \\
DDATA  & device & On each poll that produced a reading. \\
NDEATH & node   & On shutdown (published by the bridge) and on connection loss
                  (published by the broker as the last will). \\
\bottomrule
\end{tabular}
\end{center}

\code{NDATA} and \code{DDEATH} are \emph{not} emitted: the bridge carries no node-level
measurements, and a device only stops reporting when its node does.

\section{Metrics}

\begin{center}
\begin{tabular}{@{}>{\ttfamily}l l l p{5.2cm}@{}}
\toprule
\normalfont\textbf{Metric} & \textbf{Type} & \textbf{Unit} & \textbf{Meaning} \\
\midrule
Contract/Version & Int64  & ---   & Contract version, node BIRTH only. \\
Power            & Double & kW    & Instantaneous power. \\
Energy           & Double & kWh   & Cumulative energy counter. \\
bdSeq            & Int64  & ---   & Session number; BIRTH and DEATH only. \\
\bottomrule
\end{tabular}
\end{center}

Units are exact and never inferred: a source value in a unit the bridge does not
recognise is rejected, not converted on a guess.

\section{Quality and freshness}

Quality travels as a \textbf{metric property} under the key \code{Quality}, not as a
value — a consumer that ignores properties still reads a number, and one that reads them
learns whether to trust it. The codes are the conventional OPC-style triple:

\begin{center}
\begin{tabular}{@{}l r p{7.4cm}@{}}
\toprule
\textbf{Quality} & \textbf{Code} & \textbf{Meaning} \\
\midrule
\textsc{good}  & 192 & Fresh, trusted value. \\
\textsc{stale} & 500 & A previously valid value that is no longer current. \\
\textsc{bad}   & 0   & Unusable: read failure or failed validation. \\
\bottomrule
\end{tabular}
\end{center}

\subsection{Timestamps are the meter's, never the bridge's}

Every published metric is stamped with the \emph{source} measurement time
(\code{ValueDate}), never with the moment the bridge published it. Freshness is therefore
computed end-to-end, and a message lost in transit reads downstream as old data rather
than as current data.

\subsection{The first message is honest}

A cold-start device BIRTH declares the tag set with \textbf{no value} and quality
\textsc{stale}. Before the first successful fetch the bridge genuinely has nothing to
report, and declaring a placeholder zero would be indistinguishable from a meter reading
zero.

\subsection{A rebirth never promotes a value}

On reconnect the bridge re-declares its tags. Re-declared readings are degraded
(\textsc{good} becomes \textsc{stale}, never the reverse) and stamped with the reading's
own \code{ValueDate}. Without this, a 45-minute outage would come back looking like a
fresh measurement.

\section{Delivery semantics}

Every edge-node message is published at \textbf{QoS~0} with \textbf{retain false}, as the
Sparkplug specification requires. Retain is the important half: a retained payload is a
value the broker replays to any future subscriber with no way to judge its age --- a
stored lie that outlives the process. Freshness is the BIRTH's job, not the broker's.

\begin{note}
QoS~0 defines no acknowledgement, so the bridge cannot and does not claim a message was
\emph{delivered}. It reports a value as published only once the transport accepted it for
transmission, and a value it could not hand over produces a per-device traced drop rather
than silence. See ADR~0010.
\end{note}

\section{Session lifecycle}

\subsection{Boot order}

The order is not negotiable, and it is why the death certificate is built before anything
connects:

\begin{enumerate}
  \item restore \code{bdSeq} from disk and open the next session;
  \item serialise the NDEATH certificate for that session;
  \item register it as the connection's last will \emph{inside} the CONNECT packet;
  \item connect;
  \item publish NBIRTH, then one DBIRTH per meter.
\end{enumerate}

Get this wrong and the broker holds no will while the node is alive: the bridge dies and
nobody is told.

\code{bdSeq} is persisted across restarts through an atomic write (write-temp, fsync,
rename), so a power loss mid-write cannot leave a corrupt session number.

\subsection{Shutdown --- both deaths fire}
\label{sec:contract-shutdown}

On \code{SIGTERM} the bridge publishes an explicit NDEATH, keeps the transport pumping
until it reaches the wire, and then \textbf{drops} the connection. It never sends a clean
MQTT DISCONNECT, because that instructs the broker to \emph{discard} the will.

The consequence is visible on the wire and worth expecting: a consumer sees
\textbf{two NDEATH messages} for one graceful stop, carrying the same \code{bdSeq} ---
the bridge's own certificate first, then the broker's will when the socket closes at
process exit. Duplicate deaths are idempotent for a consumer that has already marked the
node down.

\begin{note}
Requiring the explicit certificate, rather than relying on the will alone, is a decision
recorded in ADR~0011. The will only fires once the broker concludes the connection is
dead; a half-open connection --- a container network torn down without a FIN, a firewall
dropping the flow silently --- stays invisible until the keep-alive elapses, which at the
configured 30~seconds means up to 45~seconds of a dead node displayed as live. That is
acceptable for an unplanned death and not for a planned stop.
\end{note}

The two mechanisms are complementary rather than alternative: the will remains the
\emph{only} mechanism for a hard death --- crash, \code{SIGKILL}, power loss.

\section{Known limitations}
\label{sec:contract-limits}

\begin{itemize}
  \item \textbf{No NCMD subscription}, so \code{Node Control/Rebirth} cannot be honoured.
        A consumer that requests a rebirth is ignored; reconnects still rebirth on their
        own.
  \item \textbf{\code{bdSeq} is fixed for the lifetime of one MQTT client}, not
        incremented per CONNECT as the specification would have it. The will is registered
        once, inside the first CONNECT, and cannot be updated afterwards; advancing the
        number on reconnect would leave the broker holding a certificate for a session
        that no longer exists, and a consumer pairing death to birth by \code{bdSeq}
        would ignore the death and keep a frozen value on screen. Matching the will is
        the honest half of that trade.
  \item \textbf{\code{bdSeq} is persisted once at boot}, not per session, and a missing or
        corrupt file restarts the numbering at the sentinel rather than refusing to start.
  \item \textbf{No report-by-exception}: an unchanged \code{ValueDate} republishes the
        same point on every poll.
  \item \textbf{Broker TLS is off} pending an upstream dependency upgrade; the broker is
        expected on a trusted local network. Traffic to the smart-me cloud is unaffected
        and remains TLS-only.
\end{itemize}

\stub{the closed meter/metric lists and the formal \code{CONTRACT\_VERSION} governance
procedure, once multi-meter support lands in a later epic.}
